\chapter{Hand Arthroplasty}

\section*{Overview}
Hand arthroplasty is joint replacement surgery of the hand, most commonly involving a finger or thumb joint damaged by arthritis, trauma, or another destructive condition. The exact approach, setting, and preparation depend on the indication, anatomy, and the patient's health.

\section*{Alternative Names}
Finger joint replacement, thumb joint replacement, small-joint arthroplasty

\section*{Principle}
Damaged joint surfaces are removed or reshaped and replaced with an implant to reduce pain and preserve useful motion. Fusion or another reconstruction may be preferable when stability is more important than motion.

\section*{Why It Is Done}
It may be considered for persistent pain, deformity, or loss of function that has not responded to medication, splinting, injections, therapy, or activity modification.

\section*{How It Is Performed}
With regional or general anesthesia, the surgeon exposes the joint, removes damaged surfaces, prepares the bone, and inserts the selected implant. Tendons and soft tissues are repaired, and the hand is splinted.

\section*{Preparation, Risks, and Recovery}
Imaging and hand-function assessment guide implant selection. Risks include infection, stiffness, nerve or tendon injury, implant loosening or wear, instability, fracture, persistent pain, and revision surgery. Therapy restores motion while protecting the repair.

\section*{Contraindications and Precautions}
Active infection, poor skin or soft-tissue condition, inadequate bone stock, uncontrolled illness, or inability to participate in rehabilitation may require postponement or a different reconstruction. The surgical team reviews hand function, nerve and blood-vessel status, smoking, and medicines affecting bleeding or healing.

\section*{When to Seek Medical Care}
Follow the procedure team's specific instructions. Seek urgent medical assessment for severe or worsening pain, uncontrolled bleeding, fever or signs of infection, trouble breathing, chest pain, fainting, new weakness or confusion, inability to urinate, or any rapidly worsening symptom. The appropriate warning signs depend on the procedure and the patient's medical condition.

\section*{References}
\begin{enumerate}
  \item MedlinePlus. Joint Replacement---Hand. U.S. National Library of Medicine; 2025.
  \item Current Medical Diagnosis and Treatment. McGraw-Hill; 2025.
\end{enumerate}

\clearpage
